\section*{文后说明：待核实问题与证据范围}
\markboth{文后核对说明}{文后核对说明}

本稿已经写出完整的方法和实验设计。表中横线表示尚无测量值，不表示零、不预设结果方向。以下为作者核对事项，不属于投稿正文。

\begin{enumerate}
\item \textbf{需要填写的实验结果。}
D/E、R0--R6、同族角色与粒度比较、跨配置多seed确认、新场景测试、机器人预训练、独立仿真/真机后训练及外部评测，尚无可填写结果。旧fixed140及1,260补集均保留开发身份。Q1九臂、探针、GT历史冻结干预以及Q2与组合的已测数字不改写为新的确认结果。

\item \textbf{本轮具体化的设计参数。}
新增1,400场确认集、三个训练seed、128轨迹的64/64、96/32、32/96分配、200\,ms/p95与80\%显存约束、2个百分点与10\%时延取舍、预训练分组划分与搜索预算、外部任务/真机速度延迟矩阵，均为本轮按完整成稿要求补齐的协议。它们不是已运行的事实。执行时应依据附录固定协议；若硬件或数据条件迫使修改，需要同时修改对应比较及其适用范围。

\item \textbf{视觉输入与真实数据资产。}
视觉状态路径指定SAM~2与FoundationPose，使用首帧目标提示及episode之前取得的物体CAD或独立参考视图。需核对DROID可用的相机标定、深度、物体参考资产和有效监督覆盖率，并填写估计误差、缺失率与成本。没有参考资产的记录不能被假定具有可用位姿；缺失输入、可用损失与数据量分别统计。所选估计器尚无本项目下的实测部署结果。

\item \textbf{最终部署实参。}
共同评测设备、相机/机械臂型号、独立标定的速度基准、实际指令周期、对象资产和各任务可用示范数需要绑定到执行记录。文中的相对速度、延迟和训练预算定义了比较关系，不表示已经采集了这些数据。真实机器人基座与两个后训练分支尚未完成，不据同步仿真或模型内计时宣称实时部署。

\item \textbf{统计与发布。}
已有主要筛选多为单训练seed，场景区间不包括训练随机性及组合选择偏差。Q1 GT历史训练具有设备切换与RNG重设履历；冻结干预衡量同一端点对错误历史的依赖。未来监督沿用原27项比较校正。投稿前还需核对CVPR2027正式模板、匿名材料与原项目发布许可；当前采用官方CVPR2026模板并设置2027会议年。
\end{enumerate}

\subsection*{旧稿中缺乏依据、或需要纠正的表述}

\begin{itemize}
\item “目标细、环境粗已经胜出”“SE(3)与flow已形成互补”“在线预测已经改善动作”“预训练或真机已验证”等结论缺少对应证据。本稿将其落实为比较对象和结果位置，不提前写成优势。
\item GT历史从824降至9/1,400是冻结策略的推理时失配效应，不是加入历史提高58.21个百分点；训练比较收益仍为5.43个百分点。该干预不证明SE(3)信息唯一或充分，也不证明视觉历史只提取SE(3)。
\item 场景流覆盖目标区域，因此P01不等于纯环境预测。已有变量替换不能单独证明粒度或二维/三维的普遍优劣。Window的历史位置读出更准而成功率更低，不能只据读出误差宣称控制收益。
\item 当前标签采用独立累计平移与左相对旋转，不能以旧候选式$\Exp(\xi^\wedge)T_t$替代。对称性损失、位姿/场景流一致性与不确定性NLL未作为既有实现加入。新Q3的标准化生成状态与旧FUTURE0原始噪声尺度分别记录。
\item OpenWAM两篇同名论文已按作者与用途区分：Wang等人的Alpha基座用于适配系统，Yu等人的可组合框架用于Q3相关工作。D参照任务网络不据此宣称继承预训练主干。CoT-VLA、Dynamic Grasping、ATM等引文的版本、作者或年份修正保留在参考文献中。
\item BASE2候选I的66.79\%来自开发池筛选，seed1为61.21\%；不能以单次筛选成绩替代最终测试。冻结夹爪分类器仅通过9/11预定回归例，其420场新评测未执行；离线改进与两例物理失败不构成未来预测改善接触的证据。
\end{itemize}
